\documentclass[parskip=half-, numbers=enddot]{scrartcl}
\usepackage[T1]{fontenc}
\usepackage[utf8]{inputenc}
\usepackage{mathtools}
\usepackage{cmbright}
\usepackage{microtype}
\usepackage{tabularray}
\UseTblrLibrary{booktabs}

\begin{document}
\begin{center}
	\bfseries
	\begin{Huge}
		Intentions pédagogiques
	\end{Huge}
	\smallbreak
	\begin{Large}
		Mathématiques - 3 GT
	\end{Large}
	\smallbreak
	\rule[10pt]{.3\linewidth}{1 pt}
\end{center}
\section{Compétences du cours de mathématiques}

\subsection{Expliciter les savoirs et les procédures}

Il s'agit d'évoquer les connaissances qui s'y rapportent et montrer
qu'on en saisit le sens, la portée. Cela se traduit par des
justifications, des explications, des erreurs à retrouver, etc.

\subsection{Appliquer des concepts et des procédures}

Au moment où l'élève apprend une procédure, il opère des raisonnements
et construit des enchaînements qui ne sont pas d'emblée des
automatismes. À son niveau, ces techniques sont parfois complexes. Il
s'agit pour lui d'acquérir des réflexes réfléchis. Il s'agira donc de
réaliser des exercices qui ressemblent fortement aux exercices déjà
rencontrés. La réalisation des exercices nécessite : une connaissance de
diverses propriétés, de règles de calcul, des procédés de constructions
des différentes figures planes, de la représentation d'objets dans
l'espace, une habileté calculatoire, la maîtrise de la
calculatrice, etc.

\subsection{Résoudre des problèmes}

La résolution de problèmes place l'élève, soit dans un contexte qui
n'est pas encore mathématisé, soit dans une situation nouvelle. Il devra
utiliser les techniques vues en classe pour les résoudre. Cela demande
d'établir des connexions entre les différents points de la matière.
Enfin, il doit pouvoir commenter, justifier les étapes de son
raisonnement et interpréter les résultats obtenus.

\section{Évaluations}

Tout au long de l'année, l'élève sera amené à évaluer sa progression et
sa compréhension des matières au travers des interrogations. Elles se
distinguent en deux catégories.

\subsection{Évaluations régulières ( Travail journalier formatif - TJF)}

Cela représente toutes les interrogations effectuées en cours de
chapitre. Elles peuvent concerner un point de matière particulier, des
exercices de drill, une vérification de l'acquisition de la théorie du chapitre, plusieurs points de matière, etc. Certaines peuvent être
programmées et d'autres non.

\subsection{Unités d'acquis d'apprentissage (Travail journalier certificatif - TFC)}

Les chapitres vus durant l'année sont regroupés au sein d'unités. À la fin de chaque unité, une évaluation certificative aura lieu. Ce sont ces évaluations ainsi que l'examen de fin d'année qui détermineront la réussite du cours de mathématiques. 
Le tableau suivant reprend les différentes unités ainsi que leur pondération.
\begin{center}
	\begin{tblr}{}
		\toprule
		Unité    &
		Chapitre & Pondération                                       \\
		\midrule
		UAA 1    & Approche graphique d'une fonction           & 15  \\
		         & Fonction du premier degré                   &     \\
		         & Systèmes de deux équations à deux inconnues &     \\
		UAA 2    & Puissance à exposants entiers               & 10  \\
		         & Inéquations                                 &     \\
		UAA 3    & Pythagore et racines                        & 10  \\
		         & Trigonométrie dans le triangles rectangle   &     \\
		UAA 4    & Polynômes                                   & 15  \\
		         & Factorisation et équations produit nul      &     \\
		         & Fractions rationnelles                      &     \\
		UAA5     & Angles et cercles                           & 10  \\
		         & Figures isométriques                        &     \\
		         & Figures semblables                          &     \\
		         & Thalès et proportions                       &     \\
		JUIN     &                                             & 40  \\
		\bottomrule
		         &                                             & 100 \\
	\end{tblr}
\end{center}

\section{Dossier d'apprentissage}

\subsection{Importance du dossier d'apprentissage}

Les dossiers d'apprentissage sont des outils précieux qui font partie
des documents officiels à conserver pour l'inspection et l'homologation.
Pour chacune des matières, le dossier (farde à devis) contient toutes
les évaluations de l'élève, lui permettant de se situer dans ses
apprentissages et permettant à ses parents de suivre le travail de leur
enfant. C'est en analysant et en corrigeant les erreurs faites lors des
évaluations que l'élève pourra progresser. Il est donc important que
chacun de ces dossiers soit bien tenu.

\subsection{Contenu et tenue du dossier d'apprentissage}

Les évaluations certificatives seront corrigées et signées avant d'être rangées dans le dossier d'apprentissage. 
Elles sont classées par ordre chronologique (du plus récent au plus ancien) et la feuille récapitulative, appelée répertoire, sera placée comme page de garde.

Les évaluations formatives quant à elles seront soit rangées dans le cours soit dans une farde à devis en fonction des désidératas du professeur. Elles seront également signées, corrigées et rangées comme les évaluations certificatives.

\section{Critère de réussite}

Le critère de réussite du cours de mathématiques est d'obtenir une note dépassant ou égalant les 50\% en cumulant les points des unités d'acquis d'apprentissage et l'examen de juin. 

Attention, un travail de remise à niveau pourra toutefois être donné fin juin aux élèves qui ont réussi l'année, mais qui ont obtenu moins de 50\% à l'examen de fin d'année.

\end{document}
